\vspace{-0.7em}
\section{Preliminary}
\label{sec:prelim}
\subsection{2D Gaussian Splatting}
2D Gaussian Splatting (2DGS)~\cite{huang20242dgs} models scenes using oriented planar disks derived from projected 3D Gaussian distributions, establishing view-consistent geometry with explicit surface normals defined as the direction of steepest density change. Each primitive is parameterized by: a center $\boldsymbol{p} \in \mathbb{R}^3$, orthogonal tangent vectors $\boldsymbol{t}_u$ and $\boldsymbol{t}_v$ defining disk orientation, and scaling factors $s_u$, $s_v$ controlling axial variances. The combined rotation and scaling transformations are represented by a covariance matrix $\boldsymbol{\Sigma}$.

Unlike volumetric rendering, 2DGS employs explicit ray-splat intersections to evaluate Gaussian contributions directly on 2D disks, enabling perspective-correct splatting while mitigating perspective distortion. The Gaussian function is defined as:
\begin{equation}
\label{eq:2dgs}
\mathcal{G}(\boldsymbol{u}) = \exp\left(-\frac{1}{2}(u^2 + v^2)\right),
\end{equation}
where $\boldsymbol{u} = (u, v)$ denotes the local UV coordinates of the ray-disk intersection.

Rendering uses front-to-back alpha blending of depth-sorted Gaussians with degeneracy handling:
\begin{equation}
\label{eq:splatting}
c(\boldsymbol{r}) = \sum_{i=1}^{N} \boldsymbol{c}_i \alpha_i \hat{\mathcal{G}}_i(\boldsymbol{u}) \prod_{j=1}^{i-1} \left(1 - \alpha_j \hat{\mathcal{G}}_j(\boldsymbol{u})\right),
\end{equation}
where $N$ is the number of overlapping Gaussians, $\alpha_i$ the opacity, and $\boldsymbol{c}_i$ the view-dependent color of the $i$-th Gaussian.
\subsection{Physically Based Deferred Shading}
In physically-based Gaussian splatting rendering, shading-rasterization order critically impacts fidelity. We adopt deferred shading to decouple this process: a geometry pass stores attributes in G-buffers, followed by a per-pixel lighting pass. Building on 3DGS-DR's pioneering deferred shading integration~\cite{ye20243dgsdr}, our method extends it to physically-based rendering using Disney BRDF~\cite{burley2012physically}. Each 2D Gaussian carries: albedo $\boldsymbol{\lambda} \in [0,1]^3$; metallic $m \in [0,1]$; roughness $r \in [0,1]$; surface normal $\boldsymbol{n} = \boldsymbol{t_u} \times \boldsymbol{t_v}$; position $\boldsymbol{x}_i$; and feature vector $k_i$ that aggregates into the feature map $K$ utilized in the second stage (Sec.~\ref{sec:method-ir}) to enhance specular effects.

We render attributes into G-buffers using alpha-blending from Eq.~\ref{eq:splatting}, replacing color $\boldsymbol{c}_i$ with $\boldsymbol{f}_i$:
\begin{equation}
\label{eq:splatting-feas}
\boldsymbol{F}=\sum_{i=1}^N \boldsymbol{f}_i \alpha_i \hat{\mathcal{G}}_i(\boldsymbol{u}(\boldsymbol{r})) \prod_{j=1}^{i-1}\left(1-\alpha_j \hat{\mathcal{G}}_j(\boldsymbol{u}(\boldsymbol{r}))\right),
\end{equation}
where $\boldsymbol{f}_i=\left[\boldsymbol{\lambda}_i, m_i, r_i, \boldsymbol{n}_i, \boldsymbol{x}_i, k_i\right]^{\top}$ contains per-Gaussian properties, yielding aggregated attributes $\boldsymbol{F}=$ $[\boldsymbol{\Lambda}, M, R, \boldsymbol{N}, \boldsymbol{D}, K]^{\top}$. 
Using aggregated G-buffer attributes, we compute outgoing radiance $L_o$ at $\boldsymbol{x}$ with normal $\boldsymbol{n}$ via the rendering equation~\cite{Kajiya_1998}:
\begin{equation}
L_o\left(\boldsymbol{\omega}_o, \boldsymbol{x}\right) = \int_{\Omega} f\left(\boldsymbol{\omega}_o, \boldsymbol{\omega}_i, \boldsymbol{x}\right) L_{\mathrm{i}}\left(\boldsymbol{\omega}_i, \boldsymbol{x}\right) (\boldsymbol{\omega}_i \cdot \boldsymbol{n})  d \boldsymbol{\omega}_i,
\label{eq:renderingEquation}
\end{equation}
where $f$ is the Disney BRDF, consisting of a diffuse term $f_d$
and a specular term $f_s$:
\begin{equation}
f_d = \frac{\boldsymbol{\Lambda}}{\pi}(1 - M)(1 - F_d),
\end{equation}
\begin{equation}
f_s\left(\boldsymbol{\omega}_i, \boldsymbol{\omega}_o, \boldsymbol{x}\right) = \frac{DFG}{4\left(\boldsymbol{\omega}_i \cdot \boldsymbol{n}\right)\left(\boldsymbol{\omega}_o \cdot \boldsymbol{n}\right)}.
\label{eq:disney_spec}
\end{equation}
Here, $D$, $F$, and $G$ represent the microsurface normal distribution, Fresnel reflectance, and geometry attenuation terms, respectively, with the diffuse Fresnel factor $F_d$ enforcing energy conservation by deducting specular-reflected energy. These terms are computed using the stored roughness $R$ and metallic $M$ values, following the energy-conserving properties of the Disney BRDF model~\cite{burley2012physically}.
%这里可以考虑补个最终辐射值的定义，如果有位置的话